\documentclass[aspectratio=169]{beamer}
\usetheme{default}
\usepackage{slidesonnet}
\title{Graph Definitions}
\begin{document}

\begin{frame}{What is a Graph?}
  \begin{itemize}
    \item A set of \textbf{vertices} (nodes)
    \item Connected by \textbf{edges}
  \end{itemize}

  \say{A graph is a mathematical structure consisting of a set of vertices, sometimes called nodes, connected by edges. Think of it like a social network where people are vertices and friendships are edges.}
\end{frame}

\begin{frame}{Euler's Theorem}
  \[ \sum_{v \in V} \deg(v) = 2|E| \]

  \say{Euler's handshaking theorem tells us that the sum of all vertex degrees equals twice the number of edges. This is because each edge contributes exactly two to the total degree count.}
\end{frame}

\begin{frame}{Dijkstra's Algorithm}
  \begin{itemize}
    \item Finds shortest paths in weighted graphs
    \item Greedy approach
  \end{itemize}

  \say{Dijkstra's algorithm is a fundamental algorithm for finding the shortest path between nodes in a weighted graph. It uses a greedy approach, always expanding the closest unvisited node.}
\end{frame}

\end{document}
